\documentclass[aspectratio=169,11pt]{beamer}
\usetheme{Madrid}
\usecolortheme{beaver}
\setbeamertemplate{navigation symbols}{}
\setbeamertemplate{footline}[frame number]
\usepackage{booktabs}
\usepackage{amsmath}

\title{Entry Regulation and the Cross-Border Cannabis Demand Premium}
\subtitle{Grant Proposal Summary}
\author{Mahyar Rezazadeh \and Babak Sahragard}
\date{\today}

\begin{document}

% ---------------- Slide 1 ----------------
\begin{frame}{Research Question \& Why It Matters}
\textbf{Research question.} Does uneven state-level cannabis legalization generate a spatial demand spillover
-- a ``border premium'' -- in legal retail markets, and does \emph{entry regulation} determine whether that
premium is absorbed by incumbent sellers' revenue or by market expansion (more, or bigger, dispensaries)?

\vspace{0.5em}
\textbf{Why it matters -- policy.}
\begin{itemize}
\item 24+ U.S.\ states have legalized recreational cannabis; state legislatures are \emph{actively debating}
license caps, moratoria, and social-equity carve-outs right now.
\item If caps redirect a demand-driven revenue gain into concentrated incumbent rents rather than broad-based
entry and tax revenue, that is a first-order, quantifiable fiscal and competition-policy concern.
\item We already estimate this cost at $\sim$\$490M in foregone border-county revenue and $\sim$\$119M in
foregone state tax revenue in a single year (CO, WA, MT), directly actionable for legislators weighing a cap.
\end{itemize}

\vspace{0.3em}
\textbf{Why it matters -- academically.} Two literatures document each half of this story separately
(cross-border cannabis demand; entry-cap market distortion) but have never been connected -- see next slide.
\end{frame}

% ---------------- Slide 2 ----------------
\begin{frame}{What We Add to the Literature}
\textbf{What is known.}
\begin{itemize}
\item \textbf{Cross-border demand is real:} Hansen, Miller \& Weber (2020, \emph{JPubE}) -- WA retailers near
Oregon lost $\sim$36\% of sales when Oregon's own market opened. Hao \& Cowan (2020, \emph{Econ.\ Inquiry});
Baggio, Chong \& Kwon (2020) confirm the mechanism in related settings.
\item \textbf{Entry caps distort market structure:} Thomas (2019) -- WA's license quota costs $\sim$18\% of
total surplus vs.\ free entry. Hollenbeck \& Uetake (2021, \emph{RAND JE}) -- WA's cap sustains meaningful
incumbent markups.
\end{itemize}

\vspace{0.4em}
\textbf{The gap.} Neither literature tests the other's variable. Border-demand papers treat the destination
market's regulation as fixed background; entry-cap papers are single-state with no external demand shock. A
systematic search (7 databases/repositories, 10+ keyword rounds, citation-tracing) found \textbf{no paper
directly testing whether entry regulation moderates the cross-border demand premium.}

\vspace{0.4em}
\textbf{Our contribution.} First 5-state design (CO, WA, OR, MT, NM) combining (i) a \emph{time-varying}
border-exposure measure that accounts for neighboring states legalizing mid-panel, (ii) an entry-regime
classification built from actual statutes rather than state labels, and (iii) direct estimation of how the
two interact -- separately for revenue per store, revenue per capita, and dispensary density.
\end{frame}

% ---------------- Slide 3 ----------------
\begin{frame}{Data Summary}
\footnotesize
\begin{table}
\centering
\begin{tabular}{lccccccc}
\toprule
State & Years & Regime & Counties & Border cty. & Interior cty. & Dispensaries \\
 & & & (range) & (range/yr) & (range/yr) & (latest yr.) \\
\midrule
Colorado   & 2015--2025 & Capped (local bans)         & 25--41 & 5--11  & 20--35 & 651 (2025) \\
Washington & 2015--2025 & Capped (state cap)          & 30--38 & 4--12  & 18--32 & 480 (2025) \\
Oregon     & 2017--2025 & Free 2017--21, Capped 2022+ & 22--30 & 1--4   & 18--27 & 777 (2025) \\
Montana    & 2022--2025 & Capped (opt-in requirement) & 23--35 & 8--11  & 15--24 & 526 (2025) \\
New Mexico & 2022--2025 & Free (statutory anti-cap)   & 29--32 & 9      & 20--23 & 696 (2025) \\
\bottomrule
\end{tabular}
\end{table}
\vspace{0.3em}
\normalsize
\begin{itemize}
\item Panel: county-year, 5 states, 2014/2015--2025 depending on each state's own retail launch (incomplete
2026 excluded); $n=964$--$1{,}283$ county-year observations depending on outcome.
\item Border exposure (70-mile threshold) is \emph{time-varying}: recomputed every year using each
neighboring state's own legalization date, not a single fixed snapshot.
\item Washington's ``capped'' status is directly confirmed: WSLCB's general license application window closed
permanently in March 2016 and has not reopened.
\end{itemize}
\end{frame}

% ---------------- Slide 4 ----------------
\begin{frame}{Estimation Strategy}
\textbf{Border premium and entry-regime interaction (main specification):}
\[
\ln(y_{ist}) = \alpha + \beta_1\,\text{Border}_{ist} + \beta_2\,\text{Capped}_{st}
+ \beta_3\,\text{Border}_{ist}\times\text{Capped}_{st} + \gamma\ln(\text{Population}_{ist})
+ \delta_t + \eta_s + \varepsilon_{ist}
\]
estimated separately for $y \in \{$revenue per store, revenue per capita, dispensary density$\}$;
county-clustered SEs. $\beta_1$ = border premium; $\beta_3$ = how capped entry changes it.

\vspace{0.3em}
\textbf{Robustness:} year-by-year event studies (Border $\times$ year); balanced panel (removes reporting
churn); county- vs.\ state$\times$year clustering; five alternative entry-cap classifications tested
head-to-head; continuous log-distance gradient in place of the binary threshold.

\vspace{0.5em}
\textbf{Revenue loss from capped entry (fiscal magnitude):}
\[
\text{Revenue gap}_s = \text{Actual revenue}_s \times \left[\frac{(1+\widehat{\text{Free premium}})}
{(1+\widehat{\text{Capped premium}})} - 1\right], \qquad
\text{Tax gap}_s = \text{Revenue gap}_s \times \tau_s
\]
where $\widehat{\text{Free premium}} = e^{\hat\beta_1}-1$, $\widehat{\text{Capped premium}} =
e^{\hat\beta_1+\hat\beta_3}-1$, and $\tau_s$ is each state's own statutory cannabis tax rate. Applied to 2025
border-county revenue in Colorado, Washington, and Montana: $\sim$\$490M revenue gap, $\sim$\$119M in foregone
state tax revenue -- reported as an illustrative order of magnitude, not a causal welfare estimate.
\end{frame}

\end{document}
